% Hafta 6 -- alıştırma gövdesi (ders notu ve ayrı alıştırma kağıdı tarafından \input edilir)
\section{Çözümlü Alıştırmalar}
Her soruyu önce kendiniz deneyin, sonra çözümü okuyun.

\begin{alistirma}
\(4x+y=12\) denklemini sağlayan kaç tane \((x,y)\) doğal sayı ikilisi vardır? \ (\(\mathbb{N}=\{0,1,2,\dots\}\))
\end{alistirma}
\begin{cozum}
\(y=12-4x\). \(y\) doğal sayı olmalı: \(12-4x\ge 0\Rightarrow 4x\le 12\Rightarrow x\le 3\).\\
\(x=0,1,2,3\ \Rightarrow\ y=12,8,4,0\). \quad İkililer: \((0,12),(1,8),(2,4),(3,0)\) \(\Rightarrow\) \(4\) tane.
\end{cozum}

\begin{alistirma}
\(\begin{cases}3x+2y=7\\ 5x-3y=-1\end{cases}\) sisteminin çözümü nedir?
\end{alistirma}
\begin{cozum}
\(y\)'yi yok edelim: 1. denklem \(\times 3\), 2. denklem \(\times 2\):
\[9x+6y=21,\qquad 10x-6y=-2\]
Taraf tarafa toplayalım: \(19x=19\Rightarrow x=1\). \quad 1. denklemde: \(3+2y=7\Rightarrow 2y=4\Rightarrow y=2\).\\
Çözüm: \((1,2)\). \quad Kontrol: \(5\cdot 1-3\cdot 2=-1\) \checkmark
\end{cozum}

\begin{alistirma}
\(\begin{cases}\dfrac{x}{3}-\dfrac{y}{2}=0\\[6pt] \dfrac{x}{2}+\dfrac{y}{4}=4\end{cases}\) sisteminin çözümü nedir?
\end{alistirma}
\begin{cozum}
Paydaları temizleyelim: 1. denklem \(\times 6\): \(2x-3y=0\); \quad 2. denklem \(\times 4\): \(2x+y=16\).\\
\(x\)'lerin katsayıları eşit; çıkaralım: \((2x+y)-(2x-3y)=16-0\Rightarrow 4y=16\Rightarrow y=4\).\\
\(2x+4=16\Rightarrow x=6\). \quad Çözüm: \((6,4)\). \quad Kontrol: \(\dfrac63-\dfrac42=0\) \checkmark, \ \(\dfrac62+\dfrac44=4\) \checkmark
\end{cozum}

\begin{alistirma}
\(4x+3y=18\) ve \(3x+4y=17\) olduğuna göre \(x+y\) ve \(x\cdot y\) kaçtır?
\end{alistirma}
\begin{cozum}
Simetrik sistem. Toplayalım: \(7x+7y=35\Rightarrow x+y=5\).\\
Çıkaralım: \(x-y=18-17=1\).\\
\(x+y=5\) ve \(x-y=1\) \(\Rightarrow\) \(2x=6\Rightarrow x=3\), \(y=2\). \quad \(x\cdot y=6\).
\end{cozum}

\begin{alistirma}
\(\begin{cases}2x+ay=3\\ 6x+9y=b\end{cases}\) sisteminin sonsuz çözümü varsa \(a+b\) kaçtır?
\end{alistirma}
\begin{cozum}
Sonsuz çözüm \(\Rightarrow\) \(\dfrac{2}{6}=\dfrac{a}{9}=\dfrac{3}{b}\). \ İlk oran \(\dfrac13\) olduğundan:
\[\frac{a}{9}=\frac13\Rightarrow a=3,\qquad \frac{3}{b}=\frac13\Rightarrow b=9,\qquad a+b=12\]
Kontrol: 1. denklemin \(3\) katı \(6x+9y=9\), yani 2. denklemin kendisi \checkmark
\end{cozum}

\begin{alistirma}
\(\begin{cases}mx-y=4\\ 2x+3y=5\end{cases}\) sisteminin tek çözümü \textbf{yoksa} \(m\) kaçtır? Bu durumda sistemin kaç çözümü vardır?
\end{alistirma}
\begin{cozum}
Tek çözüm yok \(\Rightarrow\) \(\dfrac{m}{2}=\dfrac{-1}{3}\Rightarrow m=-\dfrac23\).\\
Sabitlerin oranı \(\dfrac45\neq-\dfrac13\) olduğundan sistemin çözümü \textbf{yoktur} (\(0\) çözüm).\\
Gerçekten: \(m=-\dfrac23\) için 1. denklem \(\times(-3)\): \(2x+3y=-12\); oysa 2. denklem \(2x+3y=5\): çelişki.
\end{cozum}

\begin{alistirma}
Bir konserde öğrenci bileti \(150\) TL, tam bilet \(250\) TL'dir. Satılan toplam \(40\) biletten \(7400\) TL elde edilmiştir. Kaç öğrenci bileti satılmıştır?
\end{alistirma}
\begin{cozum}
\(x\): öğrenci bileti, \(y\): tam bilet sayısı. \quad \(x+y=40\), \ \(150x+250y=7400\Rightarrow 3x+5y=148\) (\(50\)'ye böldük).\\
\(x=40-y\) yerine koyalım: \(3(40-y)+5y=148\Rightarrow 120+2y=148\Rightarrow y=14\), \ \(x=26\).\\
Kontrol: \(150\cdot 26+250\cdot 14=3900+3500=7400\) \checkmark \quad Cevap: \(26\) öğrenci bileti.
\end{cozum}

\begin{alistirma}
İki basamaklı \(\ab{ab}\) sayısının rakamları toplamı \(12\)'dir ve \(\ab{ab}-\ab{ba}=36\)'dır. \(\ab{ab}\) kaçtır?
\end{alistirma}
\begin{cozum}
\(\ab{ab}-\ab{ba}=(10a+b)-(10b+a)=9(a-b)=36\Rightarrow a-b=4\).\\
\(a+b=12\) ile toplayalım: \(2a=16\Rightarrow a=8\), \ \(b=4\). \quad \(\ab{ab}=84\). \quad Kontrol: \(84-48=36\) \checkmark
\end{cozum}

\kendinizoncesi
\section{Alıştırmalar (Kendiniz Çözün)}

\begin{enumerate}[label=\textbf{\arabic*.},leftmargin=*,itemsep=6pt]
\item \((3,k)\) ikilisi \(2x-y=4\) denklemini sağlıyorsa \(k\) kaçtır?
\item \((-2,1)\), \((2,1)\), \((1,0)\), \((4,-1)\) ikililerinden hangileri \(x+3y=1\) denklemini sağlar?
\item \(x+y=15\) ve \(x-y=7\) ise \(x\) ve \(y\) kaçtır?
\item \(\begin{cases}2x+3y=8\\ 4x-3y=-2\end{cases}\) sisteminin çözümü nedir?
\item \(y=3x-5\) ve \(2x+y=10\) ise \(x\cdot y\) kaçtır?
\item \(\begin{cases}0{,}3x+0{,}2y=1{,}3\\ 0{,}1x-0{,}4y=-0{,}5\end{cases}\) sisteminin çözümü nedir?
\item \(x+y=10\) ve \(x-y=4\) ise \(x\cdot y\) kaçtır?
\item \(\begin{cases}(k+1)x+2y=5\\ 6x+4y=7\end{cases}\) sisteminin çözümü yoksa \(k\) kaçtır?
\item İki sayının toplamı \(63\)'tür. Sayılardan biri diğerinin \(2\) katıdır. Sayıların farkı kaçtır?
\item Bir kasada yalnızca \(10\) TL'lik ve \(20\) TL'lik banknotlardan toplam \(15\) banknot ve \(210\) TL vardır. Kaç tane \(20\) TL'lik banknot vardır?
\item Bir çiftlikteki tavuk ve tavşanların toplam \(26\) başı ve \(76\) bacağı vardır. Kaç tavuk vardır?
\item \(30\) soruluk bir testte doğru cevap \(+4\), yanlış cevap \(-1\) puandır. Bütün soruları cevaplayan bir öğrenci \(80\) puan almıştır. Kaç soruyu doğru cevaplamıştır?
\end{enumerate}

\subsection*{Cevap Anahtarı}
{\small
\begin{multicols}{2}
\begin{enumerate}[label=\textbf{\arabic*.},leftmargin=*,itemsep=1pt]
\item \(6-k=4\Rightarrow k=2\).
\item \((-2,1)\), \((1,0)\), \((4,-1)\) sağlar; \((2,1)\): \(5\neq 1\).
\item Topla: \(2x=22\Rightarrow x=11\), \ \(y=4\).
\item Topla: \(6x=6\Rightarrow x=1\); \ \(y=2\). \ \((1,2)\).
\item \(2x+3x-5=10\Rightarrow x=3\), \(y=4\); \ \(x\cdot y=12\).
\item \(\times 10\): \(3x+2y=13\), \(x=4y-5\) \(\Rightarrow\) \(14y=28\), \(y=2\), \(x=3\). \ \((3,2)\).
\item \(x=7\), \(y=3\) \(\Rightarrow\) \(x\cdot y=21\).
\item \(\dfrac{k+1}{6}=\dfrac24\Rightarrow k=2\); \ \(\dfrac57\neq\dfrac12\) \checkmark
\item \(x=2y\), \(3y=63\Rightarrow y=21\), \(x=42\); fark \(21\).
\item \(x+2y=21\), \(x+y=15\) \(\Rightarrow\) \(y=6\).
\item \(t+2r=38\), \(t+r=26\) \(\Rightarrow\) \(r=12\), \(t=14\).
\item \(d+y=30\), \(4d-y=80\) \(\Rightarrow\) \(5d=110\), \(d=22\).
\end{enumerate}
\end{multicols}}
